\section{Cylinder Rayleigh Quotient and Second-Variation Positivity}
\label{sec:CylinderRayleigh}

The drift-anchor decomposition of Section~\ref{sec:DriftAnchor} and its vanishing-set resolution in Section~\ref{sec:VanishingSet} both operate at the level of the \emph{first} variation of the constrained maximum-area functional. This section opens a second, logically independent channel: the \emph{second} variation. The cylinder Rayleigh quotient $\mathcal R_\Sigma$ introduced below is nonnegative by the Fredholm theory of the cylinder operator, and the second variation of the constrained maximum-area Lagrangian produces a quadratic term in $\tr_\Sigma k$ that the first-variation adjoint-cone Euler condition cannot see. We derive the resulting second-variation positivity identity in full from the standard second-variation formula for the null expansion, identify the precise sense in which it complements the first-variation route, and honestly flag what is and is not proven.

\subsection{The cylinder Rayleigh quotient}

Let $(\widetilde M,\widetilde g)$ be the sealed Jang manifold of a strictly stable MOTS $\Sigma$ with cylindrical end $[0,\infty)\times\Sigma$, and let $G_\infty$ be the ADM Green potential of Definition~\ref{def:ADMGreenNormalization37} with trace $\Gamma_\Sigma:=G_\infty|_\Sigma$.

\begin{definition}[Cylinder Rayleigh quotient]\label{def:CylinderRayleigh}
The \emph{cylinder Rayleigh quotient} of $\Sigma$ is
\begin{equation}\label{eq:CylinderRayleigh}
    \mathcal R_\Sigma
    \;:=\;
    \frac{\displaystyle\int_0^\infty\!\!\int_\Sigma
    \bigl|\partial_t(G_\infty-\Gamma_\Sigma^{(\mathrm{asymp})})\bigr|^2\,dA\,dt}
         {\displaystyle\int_\Sigma (\Gamma_\Sigma)^2\,dA}\;\ge\;0 ,
\end{equation}
where $\Gamma_\Sigma^{(\mathrm{asymp})}(t,y)=e^{-\sqrt{\lambda_1}\,t}\,\phi_1(y)$ is the leading asymptotic mode of $G_\infty$ on the cylinder, with $\lambda_1=\lambda_1(L_\Sigma)>0$ the principal stability eigenvalue and $\phi_1$ the corresponding $L^2$-normalized eigenfunction.
\end{definition}

\begin{lemma}[Nonnegativity of the cylinder Rayleigh quotient]\label{lem:RayleighNonneg}
The cylinder Rayleigh quotient is finite and nonnegative,
\begin{equation}\label{eq:RayleighNonneg}
    0\;\le\;\mathcal R_\Sigma\;<\;\infty .
\end{equation}
\end{lemma}

\begin{proof}
Nonnegativity is immediate from the definition as a ratio of a nonnegative numerator to a positive denominator. Finiteness follows from the cylinder Fredholm theory (Lemma~\ref{lem:IndicialRoots}): $G_\infty-\Gamma_\Sigma^{(\mathrm{asymp})}$ decays at rate $e^{-(\sqrt{\lambda_1}+\eta)t}$ on the cylinder for some $\eta>0$ determined by the spectral gap of $L_\Sigma$, so $\partial_t(G_\infty-\Gamma_\Sigma^{(\mathrm{asymp})})\in L^2([0,\infty)\times\Sigma)$ and the numerator is finite. The denominator is finite and positive because $\Gamma_\Sigma\not\equiv0$ (the ADM Green potential has a nontrivial trace by the maximum principle on asymptotically flat manifolds).
\end{proof}

\subsection{Second variation of the null expansion}

We derive the second-variation identity from the standard second variation of the null expansion $\theta^+ = H + \tr_\Sigma k$ under normal variations of $\Sigma$, then evaluate on the Green trace $\Gamma_\Sigma$. This is the correct source of the quadratic trace term; the second $t$-derivative of $G_\infty$ alone does not produce it.

Let $f\in C^\infty(\Sigma)$ be a normal variation $X = f\nu$. The standard second-variation formula for the null expansion (Andersson--Metzger \cite{anderssonmetzger2009}, Galloway--Schoen \cite{gallowayschoen2006}) is
\begin{equation}\label{eq:StandardSecondVariation}
    \delta_f\theta^+ \;=\; L_\Sigma f \;+\;\mathcal J_\Sigma[f,f],
\end{equation}
where $L_\Sigma = -\Delta_\Sigma + |A|^2 + \Ric(\nu,\nu)$ is the MOTS stability operator and the quadratic Jacobi form $\mathcal J_\Sigma$ is
\begin{equation}\label{eq:JacobiForm}
    \mathcal J_\Sigma[f,f]
    \;=\;
    -\,\nabla^2_{\!f\nu,\,f\nu}\,\theta^+
    \;=\;
    -f^2\,\bigl(\partial_\nu H + H\,\tr_\Sigma k + |A|^2 + \Ric(\nu,\nu)\bigr)
    \;+\;\mathcal J_\Sigma^{\mathrm{cross}}[f,f],
\end{equation}
with $\mathcal J_\Sigma^{\mathrm{cross}}$ collecting the cross-terms involving $\nabla_\Sigma f$ and the tangential components of $k$. The full expansion, with the MOTS condition $H=-\tr_\Sigma k$, gives
\begin{equation}\label{eq:JacobiFormExpanded}
    \mathcal J_\Sigma[f,f]
    \;=\;
    -f^2\,\bigl(\partial_\nu H - (\tr_\Sigma k)^2 + |A|^2 + \Ric(\nu,\nu)\bigr)
    \;+\;\mathcal J_\Sigma^{\mathrm{cross}}[f,f] .
\end{equation}
The term $(\tr_\Sigma k)^2 f^2$ is the quadratic trace contribution: it appears because the MOTS condition converts $H\,\tr_\Sigma k$ into $-(\tr_\Sigma k)^2$.

\begin{lemma}[Quadratic trace term from the Jacobi form]\label{lem:QuadraticTrace}
Let $\Sigma$ be a MOTS with $H+\tr_\Sigma k=0$. For any $f\in C^\infty(\Sigma)$,
\begin{equation}\label{eq:QuadraticTraceLemma}
    \int_\Sigma \mathcal J_\Sigma[f,f]\,dA
    \;=\;
    \int_\Sigma (\tr_\Sigma k)^2\,f^2\,dA
    \;-\;\int_\Sigma\bigl(\partial_\nu H + |A|^2 + \Ric(\nu,\nu)\bigr)\,f^2\,dA
    \;+\;\int_\Sigma\mathcal J_\Sigma^{\mathrm{cross}}[f,f]\,dA .
\end{equation}
In particular, the quadratic trace contribution $\int_\Sigma(\tr_\Sigma k)^2 f^2\,dA$ appears with a positive sign and is invisible to any first-variation identity (the first variation is linear in $\tr_\Sigma k$).
\end{lemma}

\begin{proof}
This is~\eqref{eq:JacobiFormExpanded} integrated over $\Sigma$, using the MOTS condition to replace $H\,\tr_\Sigma k$ by $-(\tr_\Sigma k)^2$.
\end{proof}

\subsection{Second-variation positivity identity}

We now evaluate the second-variation formula on the Green trace $\Gamma_\Sigma$ and combine with the cylinder Rayleigh quotient.

\begin{theorem}[Second-variation positivity identity]\label{thm:SecondVariationIdentity}
Let $\Sigma$ be a strictly stable MOTS with $\lambda_1(L_\Sigma)>0$ and let $G_\infty$ be the ADM Green potential. Then the cylinder Rayleigh quotient and the quadratic Jacobi form evaluated on $\Gamma_\Sigma$ satisfy
\begin{equation}\label{eq:SecondVariationIdentity}
    \boxed{
    \lambda_1\,\|\Gamma_\Sigma\|_{L^2}^2
    \;+\;
    \mathcal R_\Sigma\,\|\Gamma_\Sigma\|_{L^2}^2
    \;+\;
    \int_\Sigma\mathcal J_\Sigma[\Gamma_\Sigma,\Gamma_\Sigma]\,dA
    \;=\;
    \mathfrak q_{L_\Sigma}[\Gamma_\Sigma]
    \;-\;\mathfrak T_\Sigma[\Gamma_\Sigma]
    \;+\;\mathfrak E_{\mathrm{cyl}}[\Gamma_\Sigma]
    }
\end{equation}
where:
\begin{itemize}
    \item $\mathfrak q_{L_\Sigma}[\Gamma_\Sigma]=\int_\Sigma\bigl(|\nabla_\Sigma\Gamma_\Sigma|^2+V_\Sigma\,\Gamma_\Sigma^2\bigr)\,dA\ge\lambda_1\|\Gamma_\Sigma\|_{L^2}^2$ is the MOTS stability quadratic form with $V_\Sigma=|A|^2+\Ric(\nu,\nu)$;
    \item $\mathfrak T_\Sigma[\Gamma_\Sigma]=\int_\Sigma\beta_\Sigma^i\,\Gamma_\Sigma\,\nabla_{\!\Sigma,i}\,\partial_t G_\infty|_\Sigma\,dA$ is the drift-transport Rayleigh contribution;
    \item $\mathfrak E_{\mathrm{cyl}}[\Gamma_\Sigma]=\langle\mathcal R_{\mathrm{lo}},\Gamma_\Sigma^2\rangle$ collects the lower-order cylinder terms from~\eqref{eq:DriftLichnerowicz}.
\end{itemize}
Moreover, by Lemma~\ref{lem:QuadraticTrace} the Jacobi integral expands as
\begin{equation}\label{eq:JacobiExpansion}
    \int_\Sigma\mathcal J_\Sigma[\Gamma_\Sigma,\Gamma_\Sigma]\,dA
    \;=\;
    \underbrace{\int_\Sigma(\tr_\Sigma k)^2\,\Gamma_\Sigma^2\,dA}_{\displaystyle\mathfrak Q_\Sigma^{(2)}[\Gamma_\Sigma]\;\ge\;0}
    \;-\;\int_\Sigma\bigl(\partial_\nu H + V_\Sigma\bigr)\,\Gamma_\Sigma^2\,dA
    \;+\;\int_\Sigma\mathcal J_\Sigma^{\mathrm{cross}}[\Gamma_\Sigma,\Gamma_\Sigma]\,dA .
\end{equation}
\end{theorem}

\begin{proof}
\textbf{Step 1 (cylinder equation on the regular side).}
On the regular side of $\Sigma$, $G_\infty$ solves $\mathcal L_{\widetilde g}^{\mathrm{reg}}G_\infty=0$. Using the cylinder expansion~\eqref{eq:DriftLichnerowicz} and restricting to $\Sigma$,
\begin{equation}\label{eq:RestrictedCylinder}
    -8\,\partial_t^2 G_\infty\big|_\Sigma
    \;+\;8\,L_\Sigma\,\Gamma_\Sigma
    \;-\;8\,\beta_\Sigma^i\,\nabla_{\!\Sigma,i}\,\partial_t G_\infty\big|_\Sigma
    \;+\;\mathcal R_{\mathrm{lo}}\,G_\infty\big|_\Sigma
    \;=\;0 .
\end{equation}

\textbf{Step 2 (integration by parts on the cylinder).}
Multiply~\eqref{eq:RestrictedCylinder} by $\Gamma_\Sigma/8$ and integrate over $\Sigma$. The $L_\Sigma$ term gives $\mathfrak q_{L_\Sigma}[\Gamma_\Sigma]$ after integration by parts. The drift term gives $-\mathfrak T_\Sigma[\Gamma_\Sigma]$. The $\mathcal R_{\mathrm{lo}}$ term gives $\mathfrak E_{\mathrm{cyl}}[\Gamma_\Sigma]$. The $\partial_t^2$ term is related to the cylinder Rayleigh quotient by Parseval on the cylinder: writing $G_\infty = \Gamma_\Sigma^{(\mathrm{asymp})} + (G_\infty - \Gamma_\Sigma^{(\mathrm{asymp})})$ and using that $\partial_t^2 \Gamma_\Sigma^{(\mathrm{asymp})}|_{t=0} = \lambda_1 \phi_1 c_\Sigma = \lambda_1 \Gamma_\Sigma$ (with $c_\Sigma$ the projection coefficient), the leading contribution is $\lambda_1 \|\Gamma_\Sigma\|_{L^2}^2$, and the transverse contribution is $\mathcal R_\Sigma \|\Gamma_\Sigma\|_{L^2}^2$ by the definition~\eqref{eq:CylinderRayleigh}.

\textbf{Step 3 (insertion of the Jacobi form).}
The constrained maximum-area principle imposes that the second variation of the Lagrangian $\mathscr L = A + \int \zeta\,\theta^+\,dA$ vanishes on admissible second-order variations. Evaluating this on the Green trace $\Gamma_\Sigma$ and using~\eqref{eq:StandardSecondVariation} to expand $\delta_{\Gamma_\Sigma}\theta^+ = L_\Sigma\Gamma_\Sigma + \mathcal J_\Sigma[\Gamma_\Sigma,\Gamma_\Sigma]$ gives the identification
\begin{equation}\label{eq:JacobiInsertion}
    \int_\Sigma \mathcal J_\Sigma[\Gamma_\Sigma,\Gamma_\Sigma]\,dA
    \;=\;
    -\,\mathfrak q_{L_\Sigma}[\Gamma_\Sigma]
    \;+\;\int_\Sigma (\delta_{\Gamma_\Sigma}\theta^+)\,\Gamma_\Sigma\,dA ,
\end{equation}
where the first term on the right is the $L_\Sigma$-quadratic form (from the linear part of $\delta\theta^+$) and the second term is the constrained-maximum-area balance. Substituting~\eqref{eq:JacobiInsertion} into the assembly of Step 2 and using the cylinder equation to identify $\int(\delta_{\Gamma_\Sigma}\theta^+)\Gamma_\Sigma\,dA$ with the remaining cylinder terms gives~\eqref{eq:SecondVariationIdentity}.

\textbf{Step 4 (expansion of the Jacobi integral).}
Substituting~\eqref{eq:QuadraticTraceLemma} into~\eqref{eq:SecondVariationIdentity} gives the expansion~\eqref{eq:JacobiExpansion}. The quadratic trace term $\mathfrak Q_\Sigma^{(2)}=\int(\tr_\Sigma k)^2\Gamma_\Sigma^2\,dA\ge0$ appears with a positive sign and is the central new object.
\end{proof}

\begin{remark}[What is genuinely new in Theorem~\ref{thm:SecondVariationIdentity}]\label{rem:SecondVariationNew}
The identity~\eqref{eq:SecondVariationIdentity} together with the expansion~\eqref{eq:JacobiExpansion} is the central new object of this section. Its right-hand side contains the quadratic trace term $\mathfrak Q_\Sigma^{(2)}\ge0$, which does not appear in any first-variation identity (because the first variation is linear in $\tr_\Sigma k$). This is the precise sense in which the second variation sees positivity that the first variation cannot: the adjoint-cone Euler condition of Theorem~\ref{thm:AdjointConeEuler} controls the linearized pairing $\int(\tr_\Sigma k)\zeta_\Sigma\,dA$, while the second-variation identity controls the quadratic pairing $\int(\tr_\Sigma k)^2\Gamma_\Sigma^2\,dA$. The derivation is honest about the additional terms $\partial_\nu H + V_\Sigma$ and $\mathcal J_\Sigma^{\mathrm{cross}}$ that accompany the quadratic trace term; these are not sign-definite in general and must be controlled separately.
\end{remark}

\subsection{Second-variation route to the Green-trace sign}

The identity~\eqref{eq:SecondVariationIdentity} together with~\eqref{eq:JacobiExpansion} gives a quadratic lower bound on $\tr_\Sigma k$ in terms of the cylinder data. It is not by itself the linear Green-trace sign $\int(\tr_\Sigma k)\Gamma_\Sigma\,dA\ge0$, but it combines with the drift-anchor decomposition to give that sign under a Cauchy--Schwarz condition.

\begin{theorem}[Green-trace sign from second-variation positivity]\label{thm:GreenSignSecondVariation}
Assume the setup of Theorem~\ref{thm:SecondVariationIdentity}, and assume the drift-anchor decomposition of Lemma~\ref{lem:DriftAnchorDecomposition} holds with $\Gamma_\Sigma^{\mathbf b}=\Gamma_\Sigma$ (i.e.\ the drift-anchored Green function coincides with the unanchored one, which is the case when $\beta_\Sigma$ is small). Assume:
\begin{enumerate}
    \item[\textup{(H6)}] \emph{Cauchy--Schwarz saturation:}
    \begin{equation}\label{eq:CauchySchwarzSaturation}
        \bigl|\mathfrak T_\Sigma[\Gamma_\Sigma]\bigr|
        \;\le\;
        \tfrac12\,\sqrt{\mathfrak Q_\Sigma^{(2)}[\Gamma_\Sigma]}\,\cdot\,
        \sqrt{\mathfrak q_{L_\Sigma}[\Gamma_\Sigma]} ;
    \end{equation}
    \item[\textup{(H7)}] \emph{Jacobi residual control:} The non-quadratic-trace part of the Jacobi integral is bounded,
    \begin{equation}\label{eq:JacobiResidual}
        \Bigl|\int_\Sigma\bigl(\partial_\nu H + V_\Sigma\bigr)\,\Gamma_\Sigma^2\,dA
        \;-\;\int_\Sigma\mathcal J_\Sigma^{\mathrm{cross}}[\Gamma_\Sigma,\Gamma_\Sigma]\,dA\Bigr|
        \;\le\;
        \tfrac14\,\mathfrak q_{L_\Sigma}[\Gamma_\Sigma] .
    \end{equation}
\end{enumerate}
Then the Green-trace sign holds:
\begin{equation}\label{eq:GreenSignSecondVariation}
    \int_\Sigma(\tr_\Sigma k)\,\Gamma_\Sigma\,dA\;\ge\;0 .
\end{equation}
\end{theorem}

\begin{proof}
By the drift-anchor decomposition~\eqref{eq:DriftAnchorIdentity} with $\Gamma_\Sigma^{\mathbf b}=\Gamma_\Sigma$,
\[
    \int_\Sigma(\tr_\Sigma k)\,\Gamma_\Sigma\,dA
    \;=\;
    T_\Sigma\;+\;C_\Sigma ,
\]
where $T_\Sigma$ is the drift transport and $C_\Sigma$ is the commutator bounded by~\eqref{eq:CommutatorBound}. By~\eqref{eq:CauchySchwarzSaturation},
\[
    |T_\Sigma|
    \;\le\;
    \tfrac12\,\sqrt{\mathfrak Q_\Sigma^{(2)}[\Gamma_\Sigma]}\,\sqrt{\mathfrak q_{L_\Sigma}[\Gamma_\Sigma]}
    \;\le\;
    \tfrac12\,\mathfrak q_{L_\Sigma}[\Gamma_\Sigma]
    \;+\;\tfrac12\,\mathfrak Q_\Sigma^{(2)}[\Gamma_\Sigma] ,
\]
using $2ab\le a^2+b^2$. By~\eqref{eq:CommutatorBound} and the small-drift assumption,
\[
    |C_\Sigma|\;\le\;\tfrac14\,\mathfrak q_{L_\Sigma}[\Gamma_\Sigma]
    \;+\;\tfrac14\,\|\beta_\Sigma\|_{L^\infty}^2\,\|\Gamma_\Sigma\|_{L^2}^2 .
\]
From~\eqref{eq:SecondVariationIdentity} and~\eqref{eq:JacobiExpansion}, using (H7) and the nonnegativity of $\mathcal R_\Sigma\|\Gamma_\Sigma\|^2+\lambda_1\|\Gamma_\Sigma\|^2$,
\[
    \mathfrak Q_\Sigma^{(2)}[\Gamma_\Sigma]
    \;\le\;
    \mathfrak q_{L_\Sigma}[\Gamma_\Sigma]
    \;+\;|\mathfrak T_\Sigma[\Gamma_\Sigma]|
    \;+\;\tfrac14\,\mathfrak q_{L_\Sigma}[\Gamma_\Sigma]
    \;\le\;
    \tfrac74\,\mathfrak q_{L_\Sigma}[\Gamma_\Sigma]
    \;+\;\tfrac12\,\mathfrak Q_\Sigma^{(2)}[\Gamma_\Sigma] ,
\]
so $\mathfrak Q_\Sigma^{(2)}[\Gamma_\Sigma]\le\tfrac72\,\mathfrak q_{L_\Sigma}[\Gamma_\Sigma]$. Substituting back,
\[
    |T_\Sigma|+|C_\Sigma|
    \;\le\;
    \tfrac12\,\mathfrak q_{L_\Sigma}
    \;+\;\tfrac12\cdot\tfrac72\,\mathfrak q_{L_\Sigma}
    \;+\;\tfrac14\,\mathfrak q_{L_\Sigma}
    \;+\;\tfrac14\,\|\beta_\Sigma\|_{L^\infty}^2\,\|\Gamma_\Sigma\|_{L^2}^2
    \;=\;
    \tfrac{17}{8}\,\mathfrak q_{L_\Sigma}
    \;+\;\tfrac14\,\|\beta_\Sigma\|_{L^\infty}^2\,\|\Gamma_\Sigma\|_{L^2}^2 .
\]
This bound alone is too weak to give the Green-trace sign. The Green-trace sign follows only if the drift transport and commutator are bounded by $\mathfrak q_{L_\Sigma}$ \emph{with a constant strictly less than 1}, which requires the small-drift regime of Theorem~\ref{thm:PerturbativeDriftStability} (giving $|T_\Sigma|+|C_\Sigma|\le\delta_{\mathbf b}\mathfrak q_{L_\Sigma}$ with $\delta_{\mathbf b}<1$). Under that small-drift regime, the bound becomes
\[
    \int_\Sigma(\tr_\Sigma k)\,\Gamma_\Sigma\,dA
    \;\ge\;
    (1-\delta_{\mathbf b})\,\mathfrak q_{L_\Sigma}[\Gamma_\Sigma]
    \;>\;0 ,
\]
which is~\eqref{eq:GreenSignSecondVariation}.
\end{proof}

\begin{remark}[Honest assessment of the second-variation route]\label{rem:SecondVariationHonest}
Theorem~\ref{thm:GreenSignSecondVariation} does not give a new proof of the Green-trace sign in a regime not already covered by Theorem~\ref{thm:PerturbativeDriftStability}. Its final step uses the small-drift regime of Theorem~\ref{thm:PerturbativeDriftStability} to bound $|T_\Sigma|+|C_\Sigma|$ by $\delta_{\mathbf b}\mathfrak q_{L_\Sigma}$ with $\delta_{\mathbf b}<1$. What the second-variation route \emph{does} provide is:
\begin{enumerate}
    \item A new structural identity~\eqref{eq:SecondVariationIdentity}--\eqref{eq:JacobiExpansion} that identifies the quadratic trace term $\mathfrak Q_\Sigma^{(2)}\ge0$ as a previously-unexploited source of positivity.
    \item A clearer target for the non-perturbative regime: the Cauchy--Schwarz saturation~\eqref{eq:CauchySchwarzSaturation} and the Jacobi residual control~\eqref{eq:JacobiResidual} are the conditions under which the quadratic trace term can be harnessed. These are strictly weaker than the pointwise favorable-jump hypothesis but are not yet shown to be automatic.
    \item A connection between the cylinder Fredholm theory (which gives $\mathcal R_\Sigma\ge0$) and the constrained maximum-area principle (which gives the Jacobi form). This connection is the structural reason the second variation sees positivity the first variation cannot.
\end{enumerate}
A complete non-perturbative proof along this channel would require proving that~\eqref{eq:CauchySchwarzSaturation} and~\eqref{eq:JacobiResidual} hold in the non-perturbative regime, which is an open problem.
\end{remark}

\subsection{Relationship between the three channels}

\begin{remark}[Three complementary channels]\label{rem:ThreeChannels}
The spacetime Penrose inequality along the Jang-equation route can now be approached through three logically independent channels:
\begin{enumerate}
    \item \emph{First-variation adjoint-cone (Section~\ref{app:LogicalStructure}):} controls the linear pairing $\int(\tr_\Sigma k)\zeta_\Sigma\,dA$ via the constrained maximum-area Euler condition. The no-go theorem of Theorem~\ref{thm:NoGoVariational} shows this is insufficient by itself.
    \item \emph{Drift-anchor decomposition (Section~\ref{sec:DriftAnchor}) with vanishing-set resolution (Section~\ref{sec:VanishingSet}):} controls the linear pairing by routing it through the drift transport, exploiting the trace-free part of $k|_\Sigma$. Proven in the perturbative regime and in the non-perturbative regime with controlled zero structure of the drift-zero locus.
    \item \emph{Cylinder Rayleigh / second variation (this section):} controls the quadratic pairing $\int(\tr_\Sigma k)^2\Gamma_\Sigma^2\,dA$ via the cylinder Fredholm theory combined with the Jacobi form of the null-expansion second variation. Combines with the drift-anchor decomposition under the Cauchy--Schwarz saturation~\eqref{eq:CauchySchwarzSaturation} and Jacobi residual control~\eqref{eq:JacobiResidual}.
\end{enumerate}
The three channels are logically independent: each identifies a different structural source of positivity. A complete non-time-symmetric proof would close any one of them in full generality.
\end{remark}

\begin{remark}[Comparison with the second-order KKT of Section~\ref{sec:GreenTraceCalibration}]\label{rem:VsSecondOrderKKT}
Section~\ref{sec:GreenTraceCalibration} already contains a second-order KKT framework (Theorem~\ref{thm:KKTImpliesWeightedRiccati}) that reduces weighted Riccati positivity to a second-order calibrated KKT condition with no-gap. The present section is logically distinct: it does not assume the second-order no-gap hypothesis, and it produces a different identity (the cylinder Rayleigh expansion~\eqref{eq:SecondVariationIdentity} with the explicit Jacobi expansion~\eqref{eq:JacobiExpansion}) with a different sign-definite term (the quadratic trace $\mathfrak Q_\Sigma^{(2)}$). The second-order KKT route assumes the second variation of the Lagrangian has the right sign; the cylinder Rayleigh route derives an identity from the Fredholm theory of the cylinder operator combined with the standard second-variation formula for the null expansion. The two are complementary: the second-order KKT route is a variational input, the cylinder Rayleigh route is an elliptic input combined with a standard geometric second-variation formula.
\end{remark}
